We built and compared three restoration strategies on identical data and a fourth, generative image-to-image system, and deployed all of them as one containerised web application. The evidence supports four conclusions. (1) A single universal DAE with a genuine bottleneck learns all three corruptions (\INPUTPSNRCORR$\to$\TONEPSNRCORR\,dB), but its bottleneck caps fidelity at $\approx$28\,dB and degrades clean and mildly corrupted images. (2) Corruption classification is easy for these synthetic corruptions (99.87\,\% test accuracy), so predicted hard routing matches oracle routing; its main benefit is the identity bypass for clean inputs, and its main risk is that the rare errors --- black borders mistaken for occluders --- become severe restoration failures. Specialists were not better than the universal model on corrupted inputs, showing that the bottleneck, not task diversity, was the limiting factor. (3) The jointly trained soft MoE is the best system (\MOEPSNRCORR\,dB / \MOESSIMCORR{} SSIM on corrupted inputs) because the gate learned to blend the identity branch in proportion to corruption severity, an input-adaptive skip that hard routing cannot represent; this only became visible after we fixed a model-selection criterion that had been dominated by near-infinite identity PSNR. Soft routing did not, however, solve unseen compound corruptions, where the universal model generalised best. (4) The style-conditioned pix2pix model uses its learned embedding meaningfully --- the same face is rendered in three distinct styles --- with dense style-2 hair as the main weakness. Future work includes gated skip connections in the experts, sequential rather than parallel expert composition for compound degradations, training on mixed and real corruptions, and perceptual metrics for sketch evaluation.
